% =====================================================================
\section{Setting}\label{sec:setting}
% =====================================================================

Throughout, $\mathbf e_1=(1,0)$ and $\mathbf e_2=(0,1)$. Integer intervals are sets of integers, $[a,b]=\{a,\dots,b\}$ and $[0,m)=\{0,\dots,m-1\}$, and $z\bmod m\in[0,m)$. For $a\in\Z^2$ the shift $T_a$ acts on configurations by $(T_a\omega)(\xi)=\omega(\xi-a)$ and on measures by $T_a\mu=\mu\circ T_a^{-1}$.

\subsection{Interactions with hard constraints and Gibbs measures}

Let $\mathfrak S$ be a finite set and $\Omega=\mathfrak S^{\Z^2}$, with the product topology and its Borel $\sigma$-algebra. For $\Lambda\subseteq\Z^2$ write $\omega_\Lambda$ for the restriction of $\omega$ to $\Lambda$, and $\mathcal F_{\Lambda}$ for the $\sigma$-algebra generated by the coordinates in $\Lambda$.

\begin{definition}[Interactions]\label{set:def-interaction}
An \emph{interaction} is a family $\Phi=(\Phi_A)$, indexed by the finite sets $A\subseteq\Z^2$, of functions $\Phi_A\colon\Omega\to[0,+\infty]$ such that $\Phi_A(\omega)$ depends only on $\omega_A$. It is \emph{translation invariant} if $\Phi_{A+a}(T_a\omega)=\Phi_A(\omega)$ for all $A$, $a$ and $\omega$, and \emph{nearest-neighbour} if $\Phi_A=0$ unless $A$ is a single site or a pair $\{\xi,\xi+\mathbf e_i\}$. A configuration $\omega$ is \emph{legal} if $\Phi_A(\omega)<\infty$ for every $A$; the set of legal configurations is $X_\Phi$. For a finite set $\Lambda$, the energy of $\omega$ in $\Lambda$ is $H_\Lambda(\omega)=\sum_{A\cap\Lambda\ne\emptyset}\Phi_A(\omega)\in[0,+\infty]$.
\end{definition}

For a nearest-neighbour interaction the sum defining $H_\Lambda$ is finite. If $\omega$ is legal and $\eta$ agrees with $\omega$ outside $\Lambda$, then $\eta$ is legal exactly when $H_\Lambda(\eta)<\infty$, because the terms not meeting $\Lambda$ are the same for $\eta$ and $\omega$. We use $e^{-\infty}=0$.

\begin{definition}[Gibbs measures]\label{set:def-gibbs}
Let $\Phi$ be a nearest-neighbour interaction and $\beta>0$. For a finite $\Lambda\subseteq\Z^2$ and a legal $\omega$, the \emph{Gibbs kernel} is the probability measure on $\Omega$
\[
\gamma_\Lambda(\eta\mid\omega)=\frac{\mathbf 1[\eta_{\Lambda^c}=\omega_{\Lambda^c}]\,e^{-\beta H_\Lambda(\eta)}}{Z_\Lambda(\omega)},\qquad Z_\Lambda(\omega)=\sum_{\eta:\ \eta_{\Lambda^c}=\omega_{\Lambda^c}}e^{-\beta H_\Lambda(\eta)} .
\]
It is well defined, since $Z_\Lambda(\omega)\ge e^{-\beta H_\Lambda(\omega)}>0$, and it is supported on legal configurations. For a bounded measurable $F$ write $(\gamma_\Lambda F)(\omega)=\sum_\eta\gamma_\Lambda(\eta\mid\omega)F(\eta)$. A \emph{Gibbs measure for $\beta\Phi$} is a probability measure $\mu$ on $\Omega$ with $\mu(X_\Phi)=1$ such that
\[
\int F\,d\mu=\int(\gamma_\Lambda F)\,d\mu
\]
for every finite $\Lambda$ and every bounded measurable $F$. It is \emph{translation invariant} if $T_a\mu=\mu$ for all $a\in\Z^2$.
\end{definition}

This is the standard convention for interactions with hard constraints: Gibbs measures live on the legal configurations, and the kernels sum over the admissible fillings of $\Lambda$ only. Kernels are never evaluated at illegal boundary conditions. Since $\gamma_\Lambda F$ is $\mathcal F_{\Lambda^c}$-measurable when $F$ is $\mathcal F_\Lambda$-measurable, the identity above is equivalent to $\E_\mu[F\mid\mathcal F_{\Lambda^c}]=\gamma_\Lambda F$ $\mu$-almost surely for bounded measurable $F$.

We shall also use specifications whose variables are not indexed by sites. Let $\mathcal V$ be a countable set of \emph{coordinates}, each with a finite set of values, and let a family of nonnegative \emph{factors} be given, each a function of finitely many coordinates, with each coordinate in finitely many factors. For a finite $S\subseteq\mathcal V$ and a configuration $p$ of all coordinates, the kernel $\gamma_S(\cdot\mid p)$ is the probability on configurations agreeing with $p$ off $S$ proportional to the product of the factors that involve a coordinate of $S$, whenever the normalization is positive. A probability measure on configurations at which every factor is positive is \emph{Gibbs} for these factors if it satisfies the identity of \cref{set:def-gibbs} for every finite $S\subseteq\mathcal V$.

\begin{lemma}[Consistency]\label{set:consistency}
In the setting of factors above, let $S\subseteq S'$ be finite sets of coordinates, and let $\mu$ be a probability measure on configurations at which every factor is positive. If $\mu$ satisfies the Gibbs identity for $S'$, it satisfies it for $S$.
\end{lemma}

\begin{proof}
Fix $p$ with positive factors. The kernel $\gamma_{S'}(\cdot\mid p)$ is proportional to the product of the factors meeting $S'$; conditioning it on the values of the coordinates in $S'\setminus S$ leaves a measure proportional to the product of the factors meeting $S$, because the remaining factors are then constant. So $\gamma_{S'}(\gamma_S F)=\gamma_{S'}F$ for bounded $F$. Hence $\int F\,d\mu=\int\gamma_{S'}F\,d\mu=\int\gamma_{S'}(\gamma_SF)\,d\mu=\int\gamma_SF\,d\mu$.
\end{proof}

For an interaction on sites, the coordinates are the fields at the sites (\cref{sec:interaction}); the kernel of a set $\Lambda$ of sites is the kernel of the set of all coordinates at the sites of $\Lambda$. By \cref{set:consistency}, a measure that is Gibbs for all finite sets of sites is Gibbs for all finite sets of coordinates, since every finite set of coordinates lies in the set of coordinates of a finite set of sites; the converse is immediate.

\subsection{Turing machines}

A deterministic Turing machine has finitely many states, among them an initial state and halting states, a finite tape alphabet with a blank symbol, and a transition function. Its tape has cells $0,1,2,\dots$, its head starts on cell $0$, and at each step it writes a symbol, changes state and moves one cell left or right. We use the standard convention that the machine never moves left from cell $0$; every machine can be converted into one with this property and the same halting behaviour. A machine \emph{halts on input $n\ge1$} if, started on the tape whose cells $0,\dots,n-1$ hold the symbol $1$ and whose other cells are blank, it enters a halting state.

We fix once and for all a deterministic Turing machine $U$ of this kind that is universal: its halting set $\mathcal H_U=\{n\ge1: U\text{ halts on }n\}$ is $\Sigma^0_1$-complete, in particular undecidable. Its transition table is part of the interaction of \cref{sec:interaction}. If $U$ halts on $n$, write $t_U(n)$ for the number of steps it makes and $s_U(n)$ for the largest cell its head visits.

\subsection{Two facts about the Ising model}\label{set:sub-ising}

Let $G=(V,E)$ be a countable graph in which every vertex has at most four neighbours, and let $J\colon E\to[0,\infty)$. The \emph{Ising model} on $G$ with couplings $J$ has spins $\sigma\in\{-1,+1\}^V$ and the factors $\exp(J_{uv}\sigma_u\sigma_v)$, one for each edge $uv$. On $\Z^2$ with all couplings equal to $J$ we speak of the Ising model with coupling $J$.

\begin{lemma}[Peierls]\label{set:peierls}
For every $J\ge0$, the Ising model on $\Z^2$ with coupling $J$ has Gibbs measures $\nu^+$ and $\nu^-$, the plus and minus states (the weak limits of the finite-volume measures with plus and minus boundary conditions), invariant under all translations of $\Z^2$ and exchanged by the global spin flip. If $e^{J}\ge3$, then $\nu^+(\sigma_0=-1)\le1/3$ and $\nu^-(\sigma_0=+1)\le1/3$; in particular $\nu^+\ne\nu^-$.
\end{lemma}

\begin{proof}
The plus state $\nu^+$ is the weak limit of the finite-volume measures with plus boundary condition on boxes increasing to $\Z^2$; by the FKG inequality this limit exists, does not depend on the sequence of boxes, and is invariant under all translations, see \cite[Chapter~3]{FriedliVelenik17}. The minus state $\nu^-$ is its image under the global spin flip. By Peierls' argument \cite{Peierls36}, in a box with plus boundary condition the event $\sigma_0=-1$ forces a closed contour of the dual lattice around the origin, made of edges between disagreeing spins; a given contour of length $k$ has probability at most $e^{-2Jk}$, because flipping the spins inside it is injective and raises the weight by $e^{2Jk}$. A contour of length $k$ around the origin crosses the positive horizontal axis within distance $k/2$ of the origin, and starting from that crossing it is a closed walk of $k$ steps with at most three choices at each step after the first, so there are at most $k\,3^{k}$ of them; $k\ge4$. Hence
\[
\nu^+(\sigma_0=-1)\le\sum_{k\ge4}k\,3^{k}e^{-2Jk}\le\sum_{k\ge4}k\,3^{-k}=\frac34-\frac23=\frac1{12}\le\frac13,
\]
using $3e^{-2J}\le1/3$ and $\sum_{k\ge1}k3^{-k}=3/4$. The bound passes to the limit, and the minus state is symmetric.
\end{proof}

\begin{lemma}[Dobrushin]\label{set:dobrushin}
Let $G$ be a countable graph in which every vertex has at most four neighbours, and let $J\colon E\to[0,\bar J]$ with $4\tanh\bar J<1$. Then the Ising model on $G$ with couplings $J$ has exactly one Gibbs measure.
\end{lemma}

\begin{proof}
The conditional law of $\sigma_u$ given the other spins is $\mathbb P(\sigma_u=+1)=\frac12(1+\tanh h)$ with $h=\sum_{v\sim u}J_{uv}\sigma_v$. Changing one neighbour $\sigma_v$ changes $h$ by $2J_{uv}$, and so changes this probability by at most $\sup_t\frac12\bigl(\tanh(t+J_{uv})-\tanh(t-J_{uv})\bigr)=\tanh J_{uv}$. Dobrushin's interdependence matrix therefore has row sums at most $4\tanh\bar J<1$, and Dobrushin's uniqueness theorem \cite{Dobrushin68}, \cite[Theorem~8.7]{Georgii11} gives exactly one Gibbs measure. Existence holds because the spin space is compact and the specification has finite range.
\end{proof}

A Gibbs measure in \cref{set:dobrushin} is unique among all Gibbs measures, translation invariant or not, and the graph may be a disjoint union of infinitely many square lattices with different couplings. Both features are used in \cref{sec:proof}.
